\section{Diagrama de Gantt del proyecto}
\label{sec:gantt}

Esta sección concreta el \textit{roadmap} de doce semanas planteado en
la propuesta original (sección 7 del documento de propuesta) en un
plan calendarizado con responsables por fase, hitos académicos y
dependencias entre fases. La fecha de inicio del proyecto es el 23 de
abril de 2026 y la fecha de entrega final del MVP es el 16 de julio
de 2026.

\subsection{Tabla de fases, fechas y responsables}

La tabla \ref{tab:gantt-fases} resume las siete fases del plan, con
sus duraciones, responsables y entregables.

\begin{table}[ht]
\centering\footnotesize
\renewcommand{\arraystretch}{1.2}
\begin{tabularx}{\linewidth}{@{}c p{3.4cm} c c c c Y@{}}
\toprule
\textbf{Fase} & \textbf{Nombre} & \textbf{Sem.} & \textbf{Inicio} &
\textbf{Fin} & \textbf{Resp.} & \textbf{Entregables} \\
\midrule
F0 & Setup y configuración base & 1   & 23 abr & 29 abr & AD &
Repo, LICENSE, README, CI/CD, esqueleto \textit{back+front} \\
\addlinespace
F1 & Entrevista individual básica & 3 & 30 abr & 20 may & AD, WP &
Sala virtual, avatar LLM, voz bidireccional, aura mínima, una
industria \\
\addlinespace
F2 & Personalización y adaptación & 2 & 21 may & 03 jun & AD &
Historial, plan de mejora del LLM \textit{Coach}, perfil de
industria, calibración IRT \\
\addlinespace
F3 & \textit{Peer mock} y comunidad & 2 & 04 jun & 17 jun & AD, MS &
WebRTC, panel del observador, comentarios anclados, intercambio
de roles \\
\addlinespace
F4 & Gamificación & 1.5  & 18 jun & 27 jun & AD &
Rangos por competencia, \textit{badges}, ligas, \textit{ranking}
semanal \\
\addlinespace
F5 & Accesibilidad, ayuda y pulido & 1.5 & 28 jun & 08 jul & WP, AD &
\textit{Tooltips} del aura, asistente LLM contextual, tour guiado,
WCAG 2.2 AA \\
\addlinespace
F6 & Documentación y demo final & 1   & 09 jul & 16 jul & Todos &
Manual de usuario, documentación técnica, video demo, presentación \\
\bottomrule
\end{tabularx}
\caption{Fases del proyecto, duraciones, responsables y entregables.
Total: 12 semanas (23 abr 2026 -- 16 jul 2026).}
\label{tab:gantt-fases}
\end{table}

\subsection{Diagrama de Gantt}

La figura \ref{fig:gantt} renderiza el plan como diagrama de Gantt
con \texttt{pgfgantt}, marcando las cuatro entregas académicas
(\textit{milestones} del rombo).

\begin{figure}[ht]
\centering
\resizebox{\linewidth}{!}{%
\begin{ganttchart}[
  hgrid,
  vgrid,
  x unit=0.55cm,
  y unit chart=0.55cm,
  y unit title=0.7cm,
  bar/.append style={fill=azulUNI!75,draw=azulUNI},
  bar height=0.55,
  group/.append style={fill=azulUNI!30,draw=azulUNI},
  milestone/.append style={fill=red!70,draw=red!70!black},
  title/.style={draw=azulUNI,fill=azulUNI!10},
  title label font=\footnotesize\bfseries,
  bar label font=\footnotesize,
  group label font=\footnotesize\bfseries,
  milestone label font=\footnotesize\itshape,
]{1}{84}
  \gantttitle{Abril}{8}
  \gantttitle{Mayo}{31}
  \gantttitle{Junio}{30}
  \gantttitle{Julio}{16} \\
  \gantttitlelist{23,...,30}{1}
  \gantttitlelist{1,...,31}{1}
  \gantttitlelist{1,...,30}{1}
  \gantttitlelist{1,...,16}{1} \\

  \ganttgroup{F0 · Setup}{1}{7} \\
  \ganttbar{Repo, CI/CD, esqueleto}{1}{7} \\

  \ganttgroup{F1 · Entrevista individual}{8}{28} \\
  \ganttbar{Sala virtual + avatar}{8}{14} \\
  \ganttbar{Voz bidireccional}{15}{21} \\
  \ganttbar{Aura mínima + industria}{22}{28} \\

  \ganttgroup{F2 · Personalización}{29}{42} \\
  \ganttbar{Historial longitudinal}{29}{35} \\
  \ganttbar{Plan de mejora + IRT}{36}{42} \\

  \ganttgroup{F3 · Peer mock}{43}{56} \\
  \ganttbar{WebRTC + señalización}{43}{49} \\
  \ganttbar{Observador + comentarios}{50}{56} \\

  \ganttgroup{F4 · Gamificación}{57}{66} \\
  \ganttbar{Rangos + badges}{57}{61} \\
  \ganttbar{Ligas + ranking}{62}{66} \\

  \ganttgroup{F5 · Accesibilidad y pulido}{67}{77} \\
  \ganttbar{WCAG 2.2 AA + earcons}{67}{71} \\
  \ganttbar{Tour + tooltips + asistente}{72}{77} \\

  \ganttgroup{F6 · Documentación}{78}{84} \\
  \ganttbar{Documentación técnica}{78}{81} \\
  \ganttbar{Video demo + presentación}{82}{84} \\

  \ganttmilestone{PC02}{13} \\
  \ganttmilestone{Demo F1}{28} \\
  \ganttmilestone{Demo F3}{56} \\
  \ganttmilestone{Entrega final}{84}
\end{ganttchart}%
}
\caption{Diagrama de Gantt del plan de desarrollo (23 abr 2026 --
16 jul 2026). Los rombos marcan los hitos académicos.}
\label{fig:gantt}
\end{figure}

\subsection{Dependencias entre fases}

La tabla \ref{tab:gantt-dependencias} consolida las dependencias
lógicas que justifican la secuenciación adoptada.

\begin{table}[ht]
\centering\footnotesize
\renewcommand{\arraystretch}{1.25}
\begin{tabularx}{\linewidth}{@{}p{3.0cm} Y@{}}
\toprule
\textbf{Dependencia} & \textbf{Justificación} \\
\midrule
F0 $\rightarrow$ F1 & El repositorio y la integración continua deben
estar operativos antes de comenzar cualquier desarrollo de
funcionalidad. \\
\addlinespace
F1 $\rightarrow$ F2 & El historial longitudinal y el plan de mejora
requieren que las sesiones ya generen y persistan métricas. \\
\addlinespace
F1 $\rightarrow$ F3 & El \textit{peer mock} extiende la entrevista
individual; ambos roles requieren que la sala virtual de F1 sea
funcional. \\
\addlinespace
F2 + F3 $\rightarrow$ F4 & La gamificación necesita sesiones con
métricas persistidas (F2) y \textit{peer mock} operativo para
soportar los retos grupales (F3). \\
\addlinespace
F1\,\dots\,F4 $\rightarrow$ F5 & El pulido de accesibilidad opera
sobre las funcionalidades ya construidas, aunque los criterios WCAG
se verifican en CI desde F1. \\
\addlinespace
F5 $\rightarrow$ F6 & La documentación técnica describe el sistema
completo y requiere que el pulido de F5 esté terminado. \\
\bottomrule
\end{tabularx}
\caption{Dependencias entre fases del plan.}
\label{tab:gantt-dependencias}
\end{table}

\subsection{Posición de la PC02 en el plan}

La entrega de la PC02 (5 de mayo de 2026) cae dentro de la fase F1.
En ese punto el sistema no dispone de código completo, lo cual es
coherente con los criterios de evaluación de la PC02, que valoran el
reporte de avance y el prototipo de fidelidad media en lugar del
producto terminado. Las fases F2 a F6 quedan planificadas para las
ocho semanas restantes hasta la entrega del MVP.
